\section{Introduction}
Mathematical knowledge forms a vast network of logical dependencies that historically remained unrecorded: traditional writing preserves final results while compressing the underlying deductions and foundations, leaving the network as a latent construct obscured by coarse-grained citations. The advent of proof assistants offers the opportunity to map these dependencies, turning mathematics into explicit machine-analyzable networks.

Beginning with de~Bruijn's Automath~\cite{debruijn1970}, and followed by Mizar~\cite{trybulec1978}, Rocq/Coq~\cite{coquand1988}, and Isabelle~\cite{paulson1994}, these systems employ a verification kernel to check every declaration and record the precise premises it invokes~\cite{avigad2024}. Building on this tradition, \module{Mathlib}, the mathematical library for Lean~4~\cite{mathlib2020,lean4}, has grown to a scale where statistical structure might emerge---and where its value lies not merely in aggregating verified truths, but in transforming scattered results into a unified, queryable dependency network~\cite{feng2026erdos}.

While the validity of each theorem is enforced deterministically by the kernel, the library's macroscopic organization---naming conventions, file hierarchies, namespace boundaries---emerges from human cognitive habit and collaborative workflows. At the same time, the proof assistant is itself a programming language: Lean's type system introduces several structural layers absent from traditional mathematics, each a compromise between implicit human understanding and the compiler's demand for explicitness. The typeclass mechanism, in particular, encodes mathematical knowledge that exists as implicit shared understanding among human practitioners (e.g., ``the integers form a ring'') as explicit \emph{instance declarations} for automated reasoning. This creates an entire stratum of nodes in the dependency graph that has no traditional counterpart (\S\ref{sec:typeclass-synth}). More broadly, Lean's type system introduces distinct such compromises at every level:
\begin{itemize}
\item \emph{Coercions.} A mathematician writes ``let $n \in \mathbb{R}$'' when $n$ is a natural number and the embedding is universally understood; Lean must register an explicit chain of coercion instances $\mathbb{N} \hookrightarrow \mathbb{Z} \hookrightarrow \mathbb{Q} \hookrightarrow \mathbb{R}$, each silently inserting dependency edges.
\item \emph{Structure inheritance.} The \texttt{extends} mechanism encodes ``every group is a monoid'' not as a logical implication but as a field-packing directive that auto-generates forgetful instances, building the algebraic hierarchy's inheritance lattice.
\item \emph{Synthesized edges.} When a mathematician cites a lemma, the citation is visible on the page; when Lean resolves \texttt{a + 0 = a}, the elaborator silently inserts references to \texttt{HAdd}, \texttt{OfNat}, and the relevant instances---$74.2\%$ of all dependency edges are invisible in the source code.
\item \emph{\texttt{to\_additive}.} A textbook proves a theorem for groups and writes ``similarly for the additive case''; Lean requires a separate declaration with a distinct proof term, auto-generated by a metaprogram that duplicates every marked multiplicative result.
\item \emph{\texttt{deriving}.} Properties like decidable equality or printable representation are meta-level common sense in traditional mathematics; Lean requires explicit instance witnesses, which \texttt{deriving} handlers generate automatically.
\item \emph{Definitional height.} Mathematicians never ask how many definitions must be unfolded to verify a statement; Lean's kernel tracks this depth, which governs compilation cost and constitutes a form of implicit coupling invisible in the dependency graph.
\end{itemize}
Each mechanism deposits its own signature in the dependency graph, and collectively, they constitute the \emph{tool-infrastructure layer}. Consequently, we hypothesize that these network properties do more than document logical dependencies. Specifically, they reveal measurable traces of how mathematics is produced, reflecting community behavior, tooling mechanisms, and even the nuances of language design.

The present paper studies these structural patterns in an attempt to quantify the divergence between inherited mathematical structure (the product) and human organizational process (the process). Our goal is to establish a baseline for human-authored formal mathematics, particularly as AI-driven formalization and automated discovery increasingly capture the attention of the mathematical community~\cite{hubert2025alphaproof,achim2025aristotle,chen2025seedprover,xin2025deepseekprover}. We focus on Mathlib on \module{Mathlib}: 308,129 declarations, 8,436,366 dependency edges, and 7,563 modules. All statistics refer to commit \texttt{534cf0b} (2 Feb 2026), extracted using \texttt{importGraph} (module layer) and premise-extraction tooling (declaration layer). Our module import graph reflects a post-\texttt{shake} state (the linter removes certain transitively redundant imports in CI). Moreover, Lean~4's module system now distinguishes \texttt{public import} from \texttt{import}~\cite{lean4modules}; our snapshot captures a transitional period where most imports remain \texttt{public}. 
The rest of the paper is organized as follows. \S\ref{sec:conceptual-framework} articulates our conceptual framework: the tension between inherited mathematical structure and human organizational process, and how Lean's type system mediates this tension. \S\ref{sec:definitions} formally defines the dependency graphs and decompositions that constitute our objects of study. \S\ref{sec:contribution} states our main findings. Detailed analyses appear in the supplementary appendices: mathematical preliminaries (\S\ref{sec:preliminaries}), the three dependency graph layers (\S\ref{sec:module-import}--\S\ref{sec:namespace-graph}), structural and cross-level properties (\S\ref{sec:structural-overview}), and conclusions with implications for practice (\S\ref{sec:conclusion}).
\paragraph{Data and code availability.}
To support reproducibility and downstream research, we release the complete extraction pipeline and graph datasets. The raw graph data---including the module import graph, the declaration dependency graph, and all namespace-level aggregations---are published as a HuggingFace dataset at \texttt{MathNetwork/MathlibGraph}\footnote{\url{https://huggingface.co/datasets/MathNetwork/MathlibGraph}}, enabling independent analysis without requiring a local Lean~4 installation. The extraction and analysis code is available on GitHub at \texttt{MathNetwork/MathNetwork}\footnote{\url{https://github.com/MathNetwork/MathNetwork}}.

\section{Previous work}
\paragraph{Software dependency networks.}
Dependency graphs arise naturally in software ecosystems: package managers such as npm, PyPI, and Maven define directed acyclic graphs in which each package declares its requirements~\cite{decan2019,labelle2004}. In the formal mathematics setting, Blanchette et al.~\cite{blanchette2015} mine structural properties of the Isabelle AFP, making it the closest methodological predecessor to our work, and Bancerek et al.~\cite{bancerek2018} document the MML's role and evolution. Studies of these ecosystems have revealed heavy-tailed degree distributions, small-world connectivity, and fragility under targeted removal of hub packages. Our analysis applies a similar methodology to a formal mathematics library---a setting in which every edge carries a stronger guarantee than ``package~$A$ requires package~$B$ at build time.'' In a type-checked library, premise-level dependencies are mechanically certified: if a constant appears in a proof term, the kernel enforces that dependency. Module-level imports, by contrast, are developer-declared and can be redundant relative to reachability ($17.5\%$ are transitively implied, measured after Mathlib's \texttt{shake} linter has already pruned certain redundant imports in CI). The graph thus records not merely a build requirement, but a semantic relationship.

\paragraph{\module{Mathlib} growth and maintenance.}
Van Doorn, Ebner, and Lewis~\cite{vandoorn2020} described the engineering practices that sustain \module{Mathlib} as a collaboratively maintained library, addressing linting, documentation, and refactoring workflows. Ringer et al.~\cite{ringer2019} provide a canonical survey of proof engineering, noting that ``there is still relatively little work describing tools and best practices''; our contribution fills part of this gap. Commelin and Topaz~\cite{commelin2023} proposed abstraction boundaries and specification-driven development as organizing principles for formal mathematics. More recently, Commelin et al.~\cite{baanen2025} documented the challenges of scaling \module{Mathlib}, including build times, namespace reorganization, and technical debt. On the tooling side, Massot's \textsc{leanblueprint} system~\cite{massot2020leanblueprint} provides infrastructure for managing large formalization projects: mathematicians manually annotate dependency relations between statements in \LaTeX{} using \verb|\uses{}| tags, producing a dependency graph that tracks formalization progress. This approach has been adopted by several major projects, including the Liquid Tensor Experiment and the ongoing formalization of Fermat's Last Theorem. Zhu et al.~\cite{zhu2025leanarchitect} complement this with \textsc{LeanArchitect}, which automatically infers dependency information from Lean source code, eliminating the manual duplication between informal and formal representations. Notably, their work includes AI integration: in a case study, the automatically extracted dependency graph guides an AI prover (Achim et al.'s Aristotle system~\cite{achim2025aristotle}) to formalize theorems in topological order. Our work differs from both in scope and purpose: \textsc{leanblueprint} and \textsc{LeanArchitect} extract dependency graphs to manage individual formalization projects; we analyze the global dependency structure of the entire \module{Mathlib} library as a network, applying tools from network science---degree distributions, community detection, robustness analysis---to characterize structural patterns that no single project would reveal. Avigad et al.~\cite{avigad2025anatomy} describe the practitioner experience of using proof assistants; Klein et al.~\cite{klein2012} address managing large-scale proofs in Isabelle, including tools for moving lemmas across modules. These works address \module{Mathlib}'s maintenance from a software engineering perspective. Our analysis complements theirs by providing a network-theoretic diagnostic of the same challenges they address qualitatively: we show, for example, that the $17.5\%$ import redundancy rate (Definition~\ref{def:redundancy-rate}, measured post-\texttt{shake}) and $0.107$ module cohesion quantify the development ergonomics overhead that maintenance practices must manage.

\paragraph{Machine learning for formal mathematics.}
The graph structure of \module{Mathlib} has been exploited in machine learning, particularly for premise selection. Yang et al.~\cite{leandojo} introduced the LeanDojo benchmark, extracting proof traces from \module{Mathlib} that implicitly define a dependency graph. Bemberek et al.~\cite{bemberek2025} constructed a heterogeneous dependency graph from this benchmark and applied relational graph convolutional networks for premise selection, achieving a $25\%$ improvement over text-only baselines. Lu et al.~\cite{leanfinder} built a directed acyclic graph of Lean~4 statements with dependency edges to support semantic search. On the knowledge-representation side, the EpiK Protocol~\cite{epik2024} constructed a knowledge graph from Lean~4 and \module{Mathlib}, applying embedding models for link prediction. Uskuplu and Moss~\cite{knowtex2025} argued that dependency graphs should become a standard feature of mathematical writing. These systems consume the dependency graph as training data or as a retrieval index. We study the graph itself---not as input to a downstream task, but as an object of independent structural interest. In this context, De Moura~\cite{demoura2026proofassistants} argues that search and premise selection in \module{Mathlib} is ``not an issue'' for experienced practitioners; our structural analysis (two-layer hub structure, community detection) could inform this debate by quantifying how much of the search space is language infrastructure versus mathematical content.

\paragraph{Complex network methodology.}
The analytical tools we employ are basic and standard in network science. Heavy-tailed network theory~\cite{barabasi1999}, small-world analysis~\cite{watts1998}, and modularity-based community detection~\cite{newman2004,blondel2008} have been applied to biological, social, and technological networks, but not, to our knowledge, to a large formal mathematics library at the multi-layer depth we pursue here. (Blanchette et al.~\cite{blanchette2015} apply mining techniques to the Isabelle AFP, and Huch and Wenzel~\cite{huch2024} analyze the AFP's theory DAG for parallel build scheduling, but neither pursues the three-layer network analysis we develop.) In a related setting, Barbosa et al.~\cite{barbosa2013} studied the network structure of three online mathematical libraries---Wikipedia (restricted to mathematics), MathWorld, and DLMF---finding large strongly connected components, the absence of clear power laws, and the effectiveness of stress centrality for navigating mathematical content. Their complex networks, however, are hyperlinked encyclopedias with informal content, while here we consider a machine-verified library in which every dependency is logically certified.



